\section{Related Work}
\label{sec:relatedwork}
% \subsection{AlphaFold2 implementations}. 
% A large body of work targets training and inference efficiency on GPUs. 
% OpenFold~\cite{ahdritz2024openfold} reimplements AlphaFold2 with a focus on trainability and memory efficiency, matching AlphaFold2 accuracy while improving training throughput; industry reports describe additional engineering that shortens large-scale training on multi-GPU clusters. 
% FastFold~\cite{cheng2024fastfold} introduces Dynamic Axial Parallelism, Duality Async Operations, and AutoChunk, reporting reductions in training time and substantial speedups for long-sequence inference; later studies further systematize GPU-cluster optimizations for AlphaFold-style workloads. 
% Uni-Fold provides an open-source platform for AlphaFold with engineering improvements to training and inference pipelines. 
% 尽管这些工作在GPU上取得了一定的加速效果，但它们并未针对CPU进行优化。

% MSA and data-pipeline acceleration. Since MSA dominates wall-clock time in many deployments, the ColabFold stack replaces JackHMMER/HHblits with MMseqs2, yielding large end-to-end speedups; subsequent GPU implementations of MMseqs2 report further gains and integrate with ColabFold. These works primarily optimize the search pipeline rather than the AlphaFold2 neural kernels themselves. 

% Operator/memory specialization and long-sequence inference. Beyond parallelism, recent efforts study memory-aware chunking and kernel choices inside the Evoformer to support longer sequences with lower memory, often building on FastFold's ideas. 

% Scope relative to this work. Existing efforts overwhelmingly target GPU execution and/or MSA acceleration. To our knowledge, there is no prior systematic study that (i) optimizes AlphaFold2 inference on many-core CPUs with matrix units, and (ii) combines INT8 quantization with OPA-style kernels and NUMA-aware execution to deliver competitive end-to-end latency. Our work fills this gap by focusing on CPU-resident kernels, layout transformations, and quantization strategies tailored to AlphaFold2’s heterogeneous computation. (We base this claim on the surveys and implementations cited above.)
% Recent years have seen a surge in research focused on optimizing the training and inference efficiency of AlphaFold2, a landmark model in protein structure prediction.[1] The majority of these efforts have centered on leveraging the parallel processing capabilities of GPUs and implementing the model more efficiently.
\subsection{Optimizations for AlphaFold2}
% Since the release of AlphaFold2, a significant body of work has been dedicated to optimizing its computational pipeline for both training and inference. The majority of these efforts have focused on leveraging massively parallel hardware, primarily GPUs. For instance, OpenFold~\cite{ahdritz2024openfold} presented a faithful reimplementation of AlphaFold2 that matches the original's accuracy while significantly improving training throughput and memory efficiency, enabling training on fewer GPUs. FastFold~\cite{cheng2024fastfold} introduced several algorithmic innovations, including Dynamic Axial Parallelism and Duality Async Operations, to report substantial reductions in both training time and end-to-end inference latency, especially for long protein sequences. Similarly, Uni-Fold~\cite{li2022uni} provided an open-source PyTorch implementation with engineering enhancements that accelerated the training process compared to the official JAX-based code. ParaFold~\cite{zhong2022parafold} is an optimized system for AlphaFold inference on heterogeneous platforms, focusing on optimizing the data processing workflow. 

% While GPU-centric optimizations are prevalent, some research has explored accelerating AlphaFold2 on CPUs. The APACE framework~\cite{park2024apace} was designed to speed up AlphaFold2 in HPC environments through CPU and GPU co-optimization, parallelizing the computationally demanding MSA and template search stages. Separately, Intel has demonstrated significant inference speedups on its Xeon CPUs by leveraging hardware features like AMX with bfloat16 precision.

% In contrast to these previous efforts, our work presents the first optimization of AlphaFold2 on a many-core CPU architecture (the 72F). While prior CPU-based work focused on general parallelization or utilized bfloat16 precision on conventional CPUs, our approach is the first to exploit the dedicated matrix computation units inherent to each core of the 72F, enabling efficient low-precision inference.

Efforts to optimize AlphaFold2 have largely focused on GPUs~\cite{ahdritz2024openfold,zhu2024scalefold,cheng2024fastfold,zhong2022parafold,li2022uni,villegas2023manyfold}. OpenFold~\cite{ahdritz2024openfold} presented a faithful reimplementation of AlphaFold2 that matches the original's accuracy while significantly improving training throughput and memory efficiency, enabling training on fewer GPUs. FastFold~\cite{cheng2024fastfold} introduced several algorithmic innovations, including Dynamic Axial Parallelism and Duality Async Operations, to report substantial reductions in both training time and end-to-end inference latency, especially for long protein sequences. Similarly, Uni-Fold~\cite{li2022uni} provided an open-source PyTorch implementation with engineering enhancements that accelerated the training process compared to the official JAX-based code. ParaFold~\cite{zhong2022parafold} is an optimized system for AlphaFold inference on heterogeneous platforms, focusing on optimizing the data processing workflow. While most work is GPU-centric, some has targeted CPUs. For instance, APACE~\cite{park2024apace} used CPU and GPU co-optimization in HPC environments. Open-Omics-AlphaFold~\cite{openomicsalphafold} accelerates AlphaFold through AMX on Intel CPUs. In contrast, our work focuses on exploiting the dedicated outer product units of a many-core CPU architecture (72F) for low-precision inference, a departure from prior CPU optimizations focused on general parallelism.
% \subsection{Model Mixed-Precision and Quantization}

% Model quantization is a recognized technique for reducing computational and memory costs by using lower-precision numerical formats for model weights and activations.[19] While general-purpose INT8 quantization for large neural networks is a mature field, its application to complex models in protein structure prediction remains limited.[20][21] Some studies have investigated post-training quantization (PTQ) for protein language models like ESMFold.[22] However, these efforts have highlighted challenges, noting that standard 8-bit quantization methods can lead to a significant drop in prediction accuracy for these specialized models.[22]
% Model quantization is a well-established technique for reducing the memory footprint and computational cost of deep learning models by using lower-precision numerical formats, such as 8-bit integers (INT8). Early methods often combined quantization with knowledge distillation to help the smaller, quantized model recover the accuracy of the larger, full-precision model~\cite{polino2018model}. However, the advent of the Transformer architecture, which is a core component of AlphaFold2's Evoformer blocks, introduced new challenges.
% The primary difficulty in quantizing Transformers lies in the high dynamic range of their activation values. Large-magnitude outliers in activations can severely degrade the accuracy of standard quantization schemes. To address this, Dettmers et al. proposed GPT3.INT8()~\cite{dettmers2022gpt3}, a mixed-precision approach that performs matrix multiplication in INT8 for most features but uses FP16 for the outlier features, effectively preserving model performance. An alternative solution, SmoothQuant~\cite{xiao2023smoothquant}, introduced an offline, mathematically equivalent transformation that smooths activation outliers by migrating the quantization difficulty from activations to weights, enabling accurate and efficient INT8 quantization for both. More recently, work has also focused on specific Transformer components, with methods like SageAttention~\cite{zhang2024sageattention} being developed to enable accurate 8-bit attention calculations specifically.
% Despite these significant advances, which have been predominantly validated on large language models (LLMs), their application to scientific computing models like AlphaFold2 remains an open challenge. To our knowledge, no prior work has successfully navigated these challenges to implement and validate a full INT8 quantization scheme for the entire AlphaFold2 model. Our work is the first to do so, developing a methodology to apply both FP16 and, more importantly, INT8 precision to the AlphaFold2 inference pipeline while maintaining high accuracy. 
% Model quantization reduces compute and memory costs using low-precision formats like 8-bit integers (INT8). Significant progress has been made for large language models (LLMs), with techniques like the mixed-precision approach in GPT3.INT8() of Dettmers et al.~\cite{dettmers2022gpt3}. An alternative solution, SmoothQuant~\cite{xiao2023smoothquant}, introduced an offline, mathematically equivalent transformation that smooths activation outliers by migrating the quantization difficulty from activations to weights, enabling accurate and efficient INT8 quantization for both. More recently, work has also focused on specific Transformer components, with methods like SageAttention~\cite{zhang2024sageattention} being developed to enable accurate 8-bit attention calculations specifically.

% A central challenge in applying these methods to AlphaFold2 is the memory consumption, which is overwhelmingly dominated by activations rather than weights. The numerical distribution of these activations is highly sensitive to the input protein sequence, which makes traditional static quantization extremely difficult.
% The logical alternative is dynamic quantization, where scaling factors are computed on-the-fly for each input tensor. However, the computational overhead of calculating these factors in real-time can easily negate the performance gains from int8 matrix multiplications, posing a significant implementation challenge for achieving actual speedups.

\subsection{Other Protein Models}
Beyond AlphaFold2, other high-accuracy models have significantly advanced the field.~\cite{xu2019distance,zhou2022tasser} RoseTTAFold~\cite{baek2021accurate}, developed concurrently, achieved comparable performance using a similar deep learning architecture. Both models build upon foundational concepts from the MSA Transformer~\cite{rao2021msa}, which first demonstrated the power of attention mechanisms for processing co-evolutionary data. A different paradigm was introduced by ESMFold~\cite{lin2022language}, which leverages a large protein language model to predict structures from a single sequence, dramatically accelerating inference by bypassing the costly MSA search.

Critically, a common architectural motif unites these leading models: the extensive use of \textit{Evoformers} for iterative information refinement. This shared foundation means that the optimization techniques developed in our work are broadly applicable. Our methods for acceleration on dedicated CPU matrix units are not limited to AlphaFold2 but can be adapted to other Evoformer-based models like RoseTTAFold and ESMFold, making the proposed techniques attractive for other protein structure prediction models on CPU-centric systems.
Therefore, MatrixFold contributes to a more scalable and practical ecosystem for biological science.